\chapter{Variables}

\section{Definition and usage}

defining a variable is completely different than using it. To use it, you have to reference the variable
simlar to C.\\

\lstinputlisting[
    language=bash,
    caption={2-bash.sh},
    captionpos=t,
    label={lst:2-bash},
    breaklines=true]
{../2-bash/2-bash.sh}

\section{Variables are temporary}

This type of variable definition would not survive longer than the session. What I mean by session
is that once you close the terminal or logout that specific user, the variables are going to be erased.
This means that defining the variables using scripts can garantee that needed variables are going to 
be available for the script itself.

\section{String interpolation available}

If you know about string interpolation, you know the dollar sign is the standard way to reference a 
variable inside a string and this mechanism also works in bash.

\lstinputlisting[
    language=bash,
    caption={2-bash.sh},
    captionpos=t,
    label={lst:2.1-bash},
    breaklines=true]
{../2-bash/2.1-bash.sh}


Also, it is required to use double quotes as that is the type of string that bash understands
has the capabilities for string interpolation. If you use single strings, the variable won't be
referenced and instead, ``\$myVariable" would be directly displayed \\

Also, you can use \textbackslash{}\texttt{"} if you want to display a double quote character in your string.

\section{Store command outputs in variables}

There is this concept known as a subshell. It essentially executes a command inside a command.
This allows us to store the output of a command in a variable.

\lstinputlisting[
    language=bash,
    caption={2-bash.sh},
    captionpos=t,
    label={lst:2.2-bash},
    breaklines=true]
{../2-bash/2.2-bash.sh}

\section{Environment variables}

There are already some predefined variables. If you wish to see them, you can use the env
command

\begin{lstlisting}[
    language=bash, 
    caption={Terminal},
    captionpos=t, 
    label={lst:2-terminal}, 
    mathescape=true,
    breaklines=true]
env
\end{lstlisting}

